\section{All odd coefficients and the dyadic metabelian quotient}
\label{sec:uniform-dyadic-normalization}

The integral measure gives the exact arithmetic image in every odd
weight. Compactness then removes the entire dyadic metabelian multiplier
while preserving the arithmetic factorization.

Use the dyadic beta multiplier $B_g(X,Y)$ and the integral
homomorphisms $\tau_{2r+1}$ of
Lemma~\ref{lem:integral-boundary-measure}. Put
\[
 C_0=D,\qquad
 C_r=D\cap\bigcap_{j=1}^{r}\ker\tau_{2j+1},\qquad
 J(C_r)=C_r\cap H,\qquad C_\infty=\bigcap_{r\ge0}C_r.
\]
Here $C_1=D_3$. The finite leading-term calculations below use
$D_5=C_2$ and $D_7=C_3$. The groups $C_r$ impose only odd
metabelian coordinates; mixed coefficients will be treated separately.

\begin{theorem}\label{thm:uniform-odd-arithmetic-image}
For every integer $r\ge1$, the coefficient
\[
 m_{2r+1}(g)=[S^{2r}T]f_g\qquad(g\in C_{r-1})
\]
is a continuous homomorphism of conjugation weight
$\lambda_2^{2r+1}$, with
\[
 m_{2r+1}(C_{r-1})=m_{2r+1}(J(C_{r-1}))
                 =2^{2r+1}\mathbf Z_2.
\]
Its kernel is $C_r$. In particular each inclusion
$C_r\subset C_{r-1}$ is proper, all these subgroups are closed and
normal in $G$, and
\[
 G=C_rH\quad(r\ge0),\qquad G=C_\infty H.
\]
Moreover
\[
 C_\infty=\{g\in D:B_g(X,Y)=1\},\qquad
 B\simeq C_\infty/
       \overline{[G,C_\infty](C_\infty\cap H)}.
\]
Here $B_g$ is its dyadic power-series component; no assertion about
other components of the adelic beta function is made.
\end{theorem}

\begin{proof}
Boundary reflection makes $q_g(e^u)$ even. Since it has constant
term zero, vanishing of $\tau_3,\ldots,\tau_{2r-1}$ is equivalent
to vanishing of the coefficients of $q_g(1+S)$ below $S^{2r}$.
On $C_{r-1}$ one therefore has
\[
 \tau_{2r+1}=\frac{(2r)!}{2}\,m_{2r+1}.
\]
The asserted additivity and weight law follow, since multiplication
by a nonzero integer on $\mathbf Z_2$ is injective.
Lemma~\ref{lem:uniform-odd-coordinate-bound} gives the upper bound
$m_{2r+1}\in2^{2r+1}\mathbf Z_2$.

Choose the arithmetic elements from the cubic normalization: $\sigma\in H\cap\ker\lambda$ with
$\rho_2^I(\sigma)=2$ and all odd-prime rational Kummer coordinates
zero, and $h\in H$ with $t=\lambda(h)$ and $t_2=3$.
Then $m_3(\sigma)$ is odd. Lemma~\ref{lem:integral-boundary-measure}
gives
\[
 \tau_{2r+1}(\sigma)\equiv m_3(\sigma)\equiv1\pmod2
 \quad\text{for every }r\ge1.
\]
For a positive odd integer $j$ define
\[
 \Delta_j(z)=hzh^{-1}z^{-t^j},\qquad
 v_r=\Delta_{2r-1}\Delta_{2r-3}\cdots\Delta_3\Delta_1(\sigma).
\]
The exponent $t^j$ lies in $\widehat{\mathbf Z}^{\times}$, and the
power is taken in the closed procyclic subgroup generated by $z$.
The operator $\Delta_1$ kills every rational Kummer coordinate,
so $v_1\in D\cap H$. Each later operator kills the indicated
odd-weight homomorphism and preserves its predecessors' kernels.
Thus $v_r\in J(C_{r-1})$.

The exact value of its next coordinate is
\begin{equation}\label{eq:uniform-odd-arithmetic-witness}
 m_{2r+1}(v_r)=
 \frac{2\prod_{j=0}^{r-1}(3^{2r+1}-3^{2j+1})}{(2r)!}
                         \tau_{2r+1}(\sigma).
\end{equation}
For $k\ge1$, one has
$v_2(3^{2k}-1)=3+v_2(k)$.
For odd $k$, factor $9^k-1$ by $9-1$ and note that the remaining
sum is odd. Repeatedly squaring then adds one to the valuation,
because every $9^k+1$ has valuation one. Consequently the numerator
product in \eqref{eq:uniform-odd-arithmetic-witness}, before its
factor two, has valuation
\[
 \sum_{k=1}^{r}(3+v_2(k))=3r+v_2(r!).
\]
The denominator satisfies
$v_2((2r)!)=r+v_2(r!)$, by separating its even and odd factors.
The value in \eqref{eq:uniform-odd-arithmetic-witness} therefore has
valuation exactly $2r+1$. The closed arithmetic image contains all
profinite powers of this element, hence the whole ideal
$2^{2r+1}\mathbf Z_2$. Together with the upper bound, this proves
the two exact image equalities and identifies the kernel.

The weight laws prove normality, and the nonzero arithmetic value
proves properness. Equality of the full and arithmetic images gives
\[
 C_{r-1}=C_rJ(C_{r-1}).
\]
Starting with $G=DH$, induction gives $G=C_rH$ for all $r$.
For any fixed $g\in G$, the sets
\[
 \{a\in H:ga^{-1}\in C_r\}
\]
are nonempty, closed and nested in the compact group $H$.
Their intersection is nonempty. This proves $G=C_\infty H$.

Finally, vanishing of all $\tau_{2r+1}$ is equivalent to $q_g=0$.
The beta differential identity then gives
$B_g^{-1}Y\partial_YB_g=0$. In characteristic zero this makes
$B_g$ independent of $Y$, and $B_g(X,1)=1$ makes it identically
one. The converse follows by differentiation.
The relative-normalization argument used for
Corollary~\ref{cor:cubic-radical-normalization} applies to the
closed normal subgroup $C_\infty$ and the factorization just
proved, giving the displayed canonical quotient for $B$.
\end{proof}

\begin{proposition}\label{prop:uniform-odd-character-kernels}
For $s\ge0$, the kernel of restriction
\[
 H^1(C_s,\mathbf F_2)^G\longrightarrow
 H^1(C_\infty,\mathbf F_2)^G
\]
is the one-dimensional space generated by
\[
 \zeta_s(g)=m_{2s+3}(g)/2^{2s+3}\pmod2.
\]
For the original normalized group $N$, the kernel of restriction
to $C_\infty$ is
\[
 \left(\bigoplus_{q\text{ odd prime}}\mathbf F_2\psi_q\right)
                     \oplus\mathbf F_2\eta_3.
\]
Each of these kernels has injective restriction to the corresponding
arithmetic subgroup.
\end{proposition}

\begin{proof}
The character $\zeta_s$ is defined and invariant by the exact image
and weight law, and it vanishes on $C_{s+1}$. Its arithmetic value
on $v_{s+1}$ is one.
Let $\chi$ be an invariant character on $C_s$ vanishing on
$C_\infty$. Continuity makes $\ker\chi$ open. Compactness and the
nested intersection imply $C_r\subseteq\ker\chi$ for some
$r>s$: otherwise the compact sets $C_r\setminus\ker\chi$ would
have nonempty intersection.

Suppose $r\ge s+2$ and $\chi$ kills $C_r$. Its restriction to
$C_{r-1}$ is zero or the unique nonzero character of the quotient
$C_{r-1}/C_r\simeq\mathbf Z_2$. But
\[
 v_r=h v_{r-1}h^{-1}v_{r-1}^{-t^{2r-1}},
 \qquad v_{r-1}\in C_{r-2}\subseteq C_s.
\]
Invariance and the odd exponent give $\chi(v_r)=0$, whereas the
nonzero quotient character has value one by
\eqref{eq:uniform-odd-arithmetic-witness}. The restriction must
therefore be zero. Descending induction shows that $\chi$ kills
$C_{s+1}$, proving the first assertion.

For a character on $N$, the same compactness and descent argument
works until $C_1=D_3$: every $v_{r-1}$ used with $r\ge2$ lies
in $D\subseteq N$. Thus the kernel on $N$ is exactly its already
computed kernel on restriction to $D_3$, namely the displayed
rational Kummer span and $\eta_3$. The converse is immediate from
$C_\infty\subseteq D_3$. The arithmetic value of $\zeta_s$ and
the previously proved independence of the rational Kummer
characters and $\eta_3$ prove the last assertion.
\end{proof}

